\documentclass[11pt,a4paper]{article}
\usepackage[T1]{fontenc}
\usepackage{lmodern,amsmath,amssymb,amsthm,mathtools,booktabs,array,microtype}
\usepackage[margin=27mm,headheight=14pt]{geometry}
\usepackage[authoryear,round]{natbib}
\usepackage{fancyhdr}
\usepackage[hidelinks]{hyperref}
\usepackage{xurl}
\pagestyle{fancy}\fancyhf{}\fancyhead[L]{\small Harmonic degrees beyond the proposed band}\fancyhead[R]{\small Candidate v1.1}\fancyfoot[C]{\thepage}
\newtheorem{theorem}{Theorem}[section]
\newtheorem{lemma}[theorem]{Lemma}
\newtheorem{proposition}[theorem]{Proposition}
\newtheorem{corollary}[theorem]{Corollary}
\theoremstyle{definition}\newtheorem{definition}[theorem]{Definition}
\theoremstyle{remark}\newtheorem{remark}[theorem]{Remark}
\newcommand{\R}{\mathbb R}\newcommand{\Z}{\mathbb Z}\newcommand{\conv}{\operatorname{conv}}
\newcommand{\hbarM}{\overline h}\newcommand{\nbar}{\overline n}\newcommand{\vbar}{\overline v}
\DeclareMathOperator{\aff}{aff}\DeclareMathOperator{\relint}{relint}
\title{Rational Laguerre counterexamples to the\ harmonic-degree band for convex mosaics}
\author{Anonymous}
\date{5 October 2026\quad\textnormal{Prospective Evidence Press candidate, version 1.1}}
\begin{document}
\maketitle
\begin{center}\small Unrefereed candidate. DOI: \url{https://doi.org/10.5281/zenodo.23167002}\end{center}
\begin{abstract}
We construct a normal, periodic, face-to-face convex mosaic in three dimensions with vertex, cell and vertex--cell incidence densities
\((V,C,I)=(373,432,3276)\). Its harmonic degree is \(I/(V+C)=468/115>4\), contradicting the upper endpoint in the Domokos--L\'angi conjecture. The construction uses a rational projected deformed cube, two stackings, and a Schlegel subdivision inserted into the Freudenthal triangulation. A separate exact verifier checks all supporting inequalities, the complete face lattices, face matching and the patch volume. We also give rational periodic weighted sites whose Laguerre diagram is combinatorially dual to this mosaic. Using established projected products of polygons, we obtain harmonic degrees tending to 16; a controlled mixing argument realizes every value in \((3,16)\). Explicit cyclic-polytope families exhibit unboundedness of each mean-degree coordinate, including in the Laguerre class. Finally, we identify the closure of periodic weighted-Delaunay harmonic degrees with asymptotic vertex--facet complexity of convex four-polytopes, obtaining a lower bound of 3 in this regular class. The Schlegel insertion method and qualitative high-degree tiling constructions are prior art; the claims here are their consequences for the proposed band, the exact certified realization and the stated refinements.
\end{abstract}
\noindent\textbf{Keywords.} Convex mosaic; harmonic degree; power diagram; Schlegel diagram; neighborly cubical polytope; flag vector; exact certificate.

\section{Problem, conventions and contribution}
A convex mosaic in \(\R^d\) is a locally finite tiling by full-dimensional compact convex polytopes with disjoint interiors. It is \emph{face-to-face} when the intersection of any two cells is empty or a face of both. It is \emph{normal} when there are positive constants \(r,R\), depending on the mosaic, such that every cell contains a ball of radius \(r\) and is contained in a ball of radius \(R\). These constants need not be uniform across a sequence of different mosaics.

We use the notation of \citet{DL}: \(\nbar\) is the mean number of cells at a vertex, and \(\vbar\) is the mean number of vertices of a cell. The supplied problem bundle interchanges these two symbols; their harmonic combination is unchanged. For a periodic mosaic, write \(V,C,I\) for vertex, cell and vertex--cell incidence counts per fundamental volume, counting orbits with their incidence multiplicities. Then
\begin{equation}\label{eq:basic}
 \nbar=\frac I V,\qquad \vbar=\frac I C,\qquad
 \hbarM=\frac{\nbar\vbar}{\nbar+\vbar}=\frac I{V+C}.
\end{equation}
The same formulas hold for the densities of the nonperiodic constructions below. Bounded cell diameters and finitely many local types make the discrepancy between cell-based and vertex-based incidence counts a boundary term of order \(O(R^{d-1})\) in a radius-\(R\) observation window.

\citet[Conjecture 1]{DL} proposed \(\hbarM\in(d,2^{d-1}]\), hence \(3<\hbarM\le4\) in dimension three. Their existence theorem supplies all values in the proposed interval; their general three-dimensional lower estimate is \(28/13\). Our counterexample satisfies all the geometric hypotheses, not merely formal incidence equations.

\paragraph{Prior-art boundary.}
The underlying polytopes are not new: neighborly cubical polytopes are due to \citet{JZ}, and the larger-complexity projected products are due to \citet{Z04}. Moreover, \citet[Section 7]{Z02} explicitly describes replacing tetrahedra in a space tiling by Schlegel diagrams, discusses high mean vertex and tile degrees, and relates tiling questions to polytope complexity and liftability. Thus neither the insertion device nor qualitative one-coordinate unboundedness should be advertised as a new invention. We provide complete count formulas, a small rational counterexample to the later conjecture, a periodic power-diagram certificate, and an explicit asymptotic comparison restricted to a class for which the converse lifting is justified.

The finite counterexample is self-contained: its four-polytope is checked directly from rational inequalities in Appendix~\ref{sec:rational}. The existence of the asymptotic family approaching 16 uses the stated published projected-product theorem. The lower-bound corollary in Section~\ref{sec:equivalence} uses the established four-polytope flag inequality. These dependencies are not computational conjectures.

\begin{center}\small
\begin{tabular}{p{.27\textwidth}p{.30\textwidth}p{.33\textwidth}}
\toprule Published input & Modification established here & Conclusion here\\\midrule
Cubical polytopes \citep{JZ} & Rational specialization, two stackings, complete finite certificate & Explicit harmonic degree $468/115>4$\\
Schlegel insertion \citep{Z02}; lifting background \citep{R99} & Strict periodic implantation and rational site payload & Weighted-Delaunay and Laguerre counterexamples\\
Projected products \citep{Z04} & Exact insertion costs and unchanged-boundary mixing & Approach to 16; every value in $(3,16)$\\
Four-polytope complexity \citep{Z02} & Vanishing port cost in one direction; capped periodic epigraph in the other & Equality of finite limit sets in Theorem~\ref{thm:equivalence}\\
\bottomrule
\end{tabular}\end{center}
The general relationship between liftings and weighted diagrams is inherited background, not a novelty claim. The comparison theorem below is restricted to the stated periodic regular class.

\begin{theorem}[Certified counterexample]\label{thm:main}
There is a normal, periodic, face-to-face convex mosaic \(\mathcal M\) of \(\R^3\) with
\[
 (V,C,I)=(373,432,3276),\qquad
 (\nbar,\vbar)=\left(\frac{3276}{373},\frac{91}{12}\right),
 \qquad \hbarM=\frac{468}{115}=4+\frac8{115}.
\]
It has a rational realization. After a rational affine change of metric, it is a periodic weighted-Delaunay subdivision. Its Laguerre dual is a normal periodic convex mosaic with the two mean degrees exchanged and the same harmonic degree. Thus the proposed upper bound fails even for periodic Laguerre mosaics.
\end{theorem}

\section{Tetrahedral insertion and exact incidence accounting}
For a four-polytope \(P\), let \(f_i(P)\) count its \(i\)-faces and let \(f_{03}(P)\) count incident vertex--facet pairs. A facet of a four-polytope is a three-polytope.

\begin{lemma}[Schlegel template]\label{lem:schlegel}
Let \(P\subset\R^4\) have a tetrahedral facet \(F\). There is a convex face-to-face subdivision of \(F\) with
\[
 a=f_0(P)\text{ vertices},\quad a-4\text{ interior vertices},\quad
 b=f_3(P)-1\text{ cells},\quad j=f_{03}(P)-4\text{ incidences}.
\]
Its boundary is precisely the unsubdivided boundary of \(F\). The subdivision has a strictly regular lifting whose boundary heights are zero. Rational input data admit rational choices throughout.
\end{lemma}
\begin{proof}
Choose a point \(z\) just beyond \(F\), and strictly beneath all other facet hyperplanes. Central projection from \(z\) to \(\aff F\) maps the boundary of \(P\), with the relative interior of \(F\) removed, onto \(F\). Each ray enters \(P\) through \(F\) and exits on one of the other facets. Consequently the projected facets cover \(F\), their interiors are disjoint, and their intersections are projected common faces. Projection is projective and finite on each of these facets, so each image is a convex polytope of the same combinatorial type.

Every vertex outside \(F\) projects strictly inside \(F\). No additional boundary face or vertex appears. Removing the outer facet removes exactly four vertex--facet incidences, giving the displayed counts.

For regularity, let \(h\le0\) be the outer facet inequality and put \(s=h(z)>0\). For a point \(p\in P\), set
\[
 t(p)=\frac{s}{s-h(p)},\qquad y(p)=z+t(p)(p-z).
\]
The map \(p\mapsto(y(p),t(p)-1)\), with \(y\) coordinatized on \(\aff F\), is a nonsingular projective transformation on a neighborhood of \(P\). Its image is a convex four-polytope with top facet \(F\times\{0\}\). All remaining facets are lower facets: over an interior point of \(F\), the segment from the entrance to the exit of the ray runs downward from height zero to the exit height. Distinct facets give distinct supporting affine functions. This is the claimed strictly regular lifting. Strict inequalities define open choices, and all operations preserve rationality when the chosen viewpoint is rational.
\end{proof}

\begin{lemma}[Insertion counts]\label{lem:insert}
Suppose a tetrahedron has a convex face-to-face subdivision with \(a\) vertices, exactly four on its boundary, \(b\) cells and \(j\) vertex--cell incidences. Replacing every tetrahedron in the Freudenthal triangulation of \(\R^3\) by an affine copy of this subdivision produces a normal periodic mosaic with
\begin{equation}\label{eq:insert}
 V=1+6(a-4),\qquad C=6b,\qquad I=6j.
\end{equation}
\end{lemma}
\begin{proof}
For \(z\in\Z^3\) and a permutation \(\pi\) of \(1,2,3\), use the tetrahedron with ordered vertices
\[
 z,\quad z+e_{\pi(1)},\quad z+e_{\pi(1)}+e_{\pi(2)},\quad z+(1,1,1).
\]
These six tetrahedra triangulate every unit cube compatibly. Because the template boundary is not subdivided, neighboring copies agree exactly on every common triangle, edge and vertex. Each unit volume has one original lattice vertex and six sets of \(a-4\) new interior vertices. Cells and their incidences are assigned to unique macro-tetrahedra, proving \eqref{eq:insert}.

There are finitely many nondegenerate cell shapes, up to translation. The minimum of their positive inradii and the maximum of their circumradii establish normality. Periodicity gives the required averages. Convexity and face matching follow from the affine images of the template and its unchanged boundary.
\end{proof}

\begin{corollary}\label{cor:polytocount}
For a four-polytope with a tetrahedral facet, the inserted mosaic has
\begin{equation}\label{eq:polytocount}
 V=6f_0-23,\quad C=6(f_3-1),\quad I=6(f_{03}-4),\quad
 \hbarM=\frac{f_{03}-4}{f_0+f_3-29/6}.
\end{equation}
\end{corollary}

\section{Two stackings and the finite counterexample}
Stacking beyond a facet means taking the convex hull with a new point beyond that facet and strictly beneath every other facet. The old facet disappears and is replaced by pyramids over its two-dimensional faces. The existence of such a point follows by moving a sufficiently small distance outward from a relative-interior point of the facet.

\begin{lemma}[A tetrahedral boundary from a cube]\label{lem:port}
If a four-polytope has a cubical facet, two stackings produce a four-polytope with a tetrahedral facet, changing \((f_0,f_3,f_{03})\) by \((2,9,38)\).
\end{lemma}
\begin{proof}
Stack beyond the cube. It has six quadrilateral faces. The removed facet has eight vertices, and the six replacement square pyramids each have five vertices. Hence the changes are
\[
 \Delta f_0=1,\qquad \Delta f_3=6-1=5,\qquad
 \Delta f_{03}=6\cdot5-8=22.
\]
Now stack beyond any new square pyramid. Its five faces are one square and four triangles. The replacement facets are one square pyramid and four tetrahedra, giving
\[
 \Delta f_0=1,\qquad \Delta f_3=5-1=4,\qquad
 \Delta f_{03}=5+4\cdot4-5=16.
\]
A new tetrahedral facet is available, and addition gives the result.
\end{proof}

Neighborly cubical four-polytopes \(P_n\), \(n\ge4\), have the graph of the \(n\)-cube and cubical facets. Their established face numbers are
\begin{equation}\label{eq:ncp}
 f_0=2^n,\quad f_1=n2^{n-1},\quad
 f_2=3(n-2)2^{n-2},\quad f_3=(n-2)2^{n-2},\quad f_{03}=8f_3.
\end{equation}
For completeness, once existence and the cube graph are given, the formulas follow from \(f_0=2^n\), \(2f_1=n2^n\), ridge double counting \(f_2=3f_3\), and Euler's relation \(f_0-f_1+f_2-f_3=0\). Existence is due to \citet{JZ}. For \(n=6\) we instead provide the full rational verification below.

Applying Lemma~\ref{lem:port}, followed by Schlegel projection and insertion, gives
\begin{align}\label{eq:nfamily}
 V_n&=6\cdot2^n-11,\nonumber\\
 C_n&=6\bigl((n-2)2^{n-2}+8\bigr),\\
 I_n&=6\bigl(8(n-2)2^{n-2}+34\bigr).\nonumber
\end{align}
At \(n=6\), the complete accounting is as follows.
\begin{center}
\begin{tabular}{lrrr}
\toprule
Object & Vertices & Facets/cells & Incidences\\
\midrule
Cubical four-polytope &64&64&512\\
After first stacking &65&69&534\\
After second stacking &66&73&550\\
Tetrahedral Schlegel patch &66&72&546\\
Periodic mosaic, per unit volume &373&432&3276\\
\bottomrule
\end{tabular}
\end{center}
The patch has 62 interior vertices; the periodic count is therefore \(1+6\cdot62=373\), not \(6\cdot66\). This boundary identification is essential. Equation~\eqref{eq:basic} gives
\[
 \hbarM=\frac{3276}{373+432}=\frac{468}{115}>4,
 \qquad \vbar=\frac{3276}{432}=\frac{91}{12}.
\]
Lemmas~\ref{lem:schlegel} and \ref{lem:insert} prove the geometric assertions in Theorem~\ref{thm:main}. Rationality follows from Appendix~\ref{sec:rational}. The regular and Laguerre assertions are proved next. The family \eqref{eq:nfamily} also gives
\[
 \hbarM_n\longrightarrow8,\qquad \nbar_n\sim2n,\qquad \vbar_n\longrightarrow8.
\]

\section{A periodic weighted-Delaunay and Laguerre realization}\label{sec:regular}
A weighted site \((p,w)\) has power distance \(\|x-p\|^2-w\). Its Laguerre cell consists of points for which this power distance is minimal among the sites. The dual weighted-Delaunay complex is obtained from lower faces of the lifted points \((p,\|p\|^2-w)\). Nonsimplicial lower faces are retained; no generic perturbation or arbitrary triangulation is applied.

\begin{lemma}[Regular implantation]\label{lem:regular}
A finite strictly regular tetrahedral subdivision with no boundary subdivision can be implanted periodically in the Freudenthal triangulation so that the resulting mosaic is a weighted-Delaunay subdivision in a Euclidean metric. Rational template coordinates and heights give rational periodic site coordinates and weights after a rational linear transformation.
\end{lemma}
\begin{proof}
Let \(J\) be the \(3\times3\) all-ones matrix and choose
\[
 Q=I-\frac5{16}J,\qquad A=I-\frac14J,\qquad A^{\mathsf T}A=Q.
\]
The eigenvalues of \(Q\) are \(1,1,1/16\), so \(q(x)=x^{\mathsf T}Qx\) is positive definite. Interpolate its lattice values affinely on the Freudenthal tetrahedra, obtaining \(L\).

We verify that \(L\) is convex and has a strict bend at each macro-face. Inside a unit cube, its simplices correspond to descending orders of the fractional coordinates. If two adjacent coordinates exchange order, the convexity gap is \(-2Q_{ij}>0\). Across a coordinate-integer plane, the entering and leaving simplices have gradient jump \(2\sum_jQ_{ij}=1/8>0\) in the normal coordinate. These are all macro-face types. Thus every bend is strict. The interpolation also satisfies
\begin{equation}\label{eq:quasiper}
 L(x+z)=L(x)+2(Qz)\cdot x+q(z),\qquad z\in\Z^3.
\end{equation}

Let \(g\) be the convex piecewise-affine template lifting, normalized to zero on the tetrahedron boundary. Transfer the same \(g\) affinely to each macro-tetrahedron, and call the resulting continuous function \(G\). It is convex inside each macro-tetrahedron, although it can have concave bends on macro-faces. Only finitely many pairs of adjacent template types occur. Consequently for all sufficiently small \(\delta>0\),
\[
 H=L+\delta G
\]
is convex on the whole space and still strictly bent at every macro-face. Within a macro-tetrahedron its maximal affine pieces are exactly the template cells. Local convexity across all faces implies global convexity by restricting to a generic line segment and observing nondecreasing successive slopes, then taking limits for nongeneric segments.

The inserted vertex set has finitely many representatives \(x_i\) modulo \(\Z^3\). Use Euclidean sites \(p_i=Ax_i\) and weights
\[
 w_i=q(x_i)-H(x_i).
\]
Because \(G\) is periodic and \eqref{eq:quasiper} holds, the weight is unchanged when \(x_i\) is translated by \(\Z^3\). All lifted site heights are \(H(x_i)\); convexity says their lower faces are precisely the constructed mosaic. A rational \(\delta\) can be chosen using finitely many strict rational inequalities.
\end{proof}

\begin{lemma}[Duality and normality]\label{lem:dual}
For the periodic weighted-Delaunay mosaics above, the Laguerre dual is a normal periodic face-to-face convex mosaic. Its density triple is \((C,V,I)\), and its mean degrees are \((\vbar,\nbar)\).
\end{lemma}
\begin{proof}
An exposed lower facet corresponds to a Laguerre vertex; an exposed lower vertex corresponds to a full-dimensional Laguerre cell. The lower hull therefore gives an incidence-reversing duality between the two polyhedral complexes, including nonsimplicial faces. Strict convexity across every cell face ensures that no listed facet is artificially split into coplanar cells.

There are finitely many site and weight types per period, the sites are relatively dense, and the weights are bounded. Comparing a distant site with a uniformly nearby site shows that every nonempty power cell is bounded, with a uniform diameter bound. All active sites correspond to full-dimensional power cells. There are finitely many such cells up to lattice translation, so their positive inradii have a positive minimum. Power cells are convex intersections of halfspaces and fit face-to-face. Dual incidences have bounded spatial displacement, hence agree in density, giving the count exchange.
\end{proof}

\paragraph{Exact weighted-site certificate.}
For the 66-vertex patch the script \path{verify_regular_realization.py} obtains
\[
 \delta=\frac1{256}.
\]
It checks all 24 directed macro-face incidences adjacent to the six tetrahedra in a unit cube. Every exact convexity gap is positive; their minimum in the chosen opposite-corner normalization is
\[
 \frac{8426635718028359287}{255182159259348442040}>0.
\]
The resulting file \texttt{weighted\_sites\_certificate.json} lists 373 rational site representatives, their rational weights, the rational lattice matrix \(A\), and all 432 dual-Delaunay cell orbits. It defines a periodic power diagram without needing an infinite search. The Laguerre mosaic has \((V,C,I)=(432,373,3276)\). This completes Theorem~\ref{thm:main}.

\section{The attainable mean-degree set is unbounded}\label{sec:degrees}
Let \(S_d\) denote attainable \((\nbar,\vbar)\) pairs in normal face-to-face convex \(d\)-mosaics with existing means.

\begin{theorem}\label{thm:degrees}
For every \(d\ge3\), \(S_d\) is unbounded in each coordinate separately. In dimension three there are periodic weighted-Delaunay mosaics with \(\vbar=4\) and \(\nbar\to\infty\), and periodic Laguerre mosaics with \(\nbar=4\) and \(\vbar\to\infty\).
\end{theorem}
\begin{proof}
Take the cyclic four-polytope \(C_4(m)\), \(m\ge5\), the convex hull of \((t,t^2,t^3,t^4)\) at \(m\) distinct real values of \(t\). It is simplicial and two-neighborly. One elementary proof of the edge property is to use the nonnegative quartic \((t-a)^2(t-b)^2\), which exposes any two chosen vertices. No five vertices are coplanar, by the Vandermonde determinant, so every facet is a tetrahedron. Thus \(f_1=\binom m2\), and ridge double counting and Euler give
\[
 a=m,\qquad b=f_3=\binom m2-m=\frac{m(m-3)}2,\qquad f_{03}=4b.
\]
Schlegel projection through any facet and insertion yield
\[
 V=6m-23,\qquad C=6(b-1),\qquad I=24(b-1).
\]
Consequently
\begin{equation}\label{eq:cyclicmeans}
 \vbar=4,\qquad \nbar=\frac{24(b-1)}{6m-23}\sim2m.
\end{equation}
Lemmas~\ref{lem:regular} and \ref{lem:dual} produce the weighted-Delaunay realization and its Laguerre dual, proving both three-dimensional statements.

For \(d=3+k\), take the Cartesian product with the unit-interval tiling in \(k\) additional coordinates. A product cell has \(2^k\) times as many vertices, and a product vertex meets \(2^k\) times as many cells. Normality, face matching and existence of the means are preserved. Hence both degrees, and the harmonic degree, are multiplied by \(2^k\).
\end{proof}

This is a direct quantitative treatment of the boundedness question in the supplied bundle. The qualitative phenomenon is already discussed in \citet[Section 7]{Z02}; we do not claim priority for unboundedness itself. The preceding proof also shows exactly why an unbounded coordinate does not force unbounded harmonic degree: in \eqref{eq:cyclicmeans}, \(\hbarM\to4\).

\begin{proposition}[The planar comparison]\label{prop:planar}
The full planar set is
\[
 S_2=\left\{(n,v):3\le n,v\le6,\quad \frac1n+\frac1v=\frac12\right\}.
\]
\end{proposition}
\begin{proof}
Euler's relation in large windows, together with edge--cell and edge--vertex incidence counting, gives \(\hbarM=2\). Every node and every polygon has degree at least three, so the displayed bounds follow.

Conversely, start from the regular hexagonal tiling. In a proportion \(p\in[0,1]\) of its hexagons, insert the center and join it to all six vertices, leaving all hexagon boundaries unchanged. Per original hexagon the densities are
\[
 V=2+p,\qquad C=1+5p,\qquad I=6+12p.
\]
As \(p\) varies, \(I/V\) increases continuously from 3 to 6 and the required curve is traced exactly. Rational proportions can be periodic. For any real proportion use the balanced column sequence \(a_j=\lfloor(j+1)p\rfloor-\lfloor jp\rfloor\). Its bounded interval discrepancy gives the same densities in expanding disks, by summation by parts. There are only the original hexagons and the six congruent triangular cell types, so the mosaic is normal.
\end{proof}

\section{Harmonic degrees approaching 16 and interval filling}\label{sec:16}
The following published input is enough; no numerical search for high-incidence cells is required.

\begin{proposition}[Projected-product input, \citet{Z04}]\label{prop:product}
For even \(m\ge4\) and integers \(r\ge3\), there is a convex four-polytope \(P_{r,m}\) with
\begin{equation}\label{eq:productcounts}
 a=m^r,\quad b=\frac{r-2}{4}m^r+r m^{r-1},\quad
 j=4(r-1)m^r.
\end{equation}
Its facets consist of \((r-2)m^r/4\) combinatorial cubes and \(r m^{r-1}\) prisms over \(m\)-gons.
\end{proposition}
The stated \(b,j\) can be checked independently from the facet types: each cube has eight vertices and each prism has \(2m\), so \(j=8((r-2)m^r/4)+2m(rm^{r-1})\). The geometric existence is the nontrivial published input. The facet count is the sum of the cube and prism counts; the vertex--facet incidence count weights these by eight and \(2m\), respectively. No other flag-vector component is used here.

\begin{theorem}\label{thm:16}
There are periodic normal weighted-Delaunay mosaics, and their Laguerre duals, with harmonic degrees tending to 16. Moreover every \(h\in(3,16)\) is attained by a normal face-to-face convex mosaic. Every rational \(h\in(3,16)\) is attained periodically.
\end{theorem}
\begin{proof}
There is a cube facet in Proposition~\ref{prop:product}. Apply the two-stack construction and insert its tetrahedral Schlegel diagram. The resulting harmonic degree is
\begin{equation}\label{eq:hproduct}
 h_{r,m}=\frac{j+34}{a+b+37/6}
 =\frac{4(r-1)m^r+34}{m^r+(r-2)m^r/4+r m^{r-1}+37/6}.
\end{equation}
The additive constants are negligible, and
\[
 \frac{j}{a+b}=\frac{16(r-1)}{r+2+4r/m}\longrightarrow16
 \quad\text{as }r,m\longrightarrow\infty.
\]
Regular implantation and duality apply exactly as before.

For low harmonic degrees, begin with the Freudenthal triangulation and perform \(k\) rounds in which every tetrahedron is subdivided into four by joining its center to its facets. Each round leaves the outer boundary of every original tetrahedron unchanged. Direct induction gives
\[
 C_k=6\cdot4^k,\qquad V_k=2\cdot4^k-1,\qquad I_k=24\cdot4^k,
 \qquad h_k=\frac{24\cdot4^k}{8\cdot4^k-1}\downarrow3.
\]
Fix \(h\in(3,16)\). Choose a low template with harmonic degree below \(h\), and a high template from \eqref{eq:hproduct} above \(h\). They have the same unsubdivided tetrahedral boundary. Place the high template in a proportion \(p\) of macro-tetrahedra and the low template in the rest. Each of \(V,C,I\) is affine in \(p\), with the original lattice-vertex contribution included once. Therefore \(I/(V+C)\) is a continuous fractional-linear function connecting the two endpoint values; its derivative has a fixed nonzero sign. It takes the value \(h\) for a unique \(p\in(0,1)\).

A periodic selection realizes every rational \(p\). All endpoint densities are rational, so solving the fractional-linear equation for a rational \(h\) gives rational \(p\). For irrational \(p\), assign the same selection to all six tetrahedra of cube-column \(z_1=j\), using \(\lfloor(j+1)p\rfloor-\lfloor jp\rfloor\). Interval discrepancy is bounded by one. Summation by parts against the cross-sectional counts of a large ball gives an \(O(R^2)\) density error, while the total counts are of order \(R^3\). Finitely many cell types ensure normality. Thus every claimed mean exists and every claimed value is attained.
\end{proof}

\begin{remark}
The interval theorem does not say that 16 is the maximum, or that 3 cannot be attained in the full convex class. The endpoints in the proof are limits. Periodic mosaics necessarily have rational mean degrees, so irrational harmonic degrees require nonperiodic constructions. The examples obtained by taking products in Theorem~\ref{thm:degrees} also contradict the upper endpoint \(2^{d-1}\) of the proposed band in every \(d\ge3\).
\end{remark}

\section{Asymptotic flag complexity and the regular-class frontier}\label{sec:equivalence}
Write
\[
 \rho(P)=\frac{f_{03}(P)}{f_0(P)+f_3(P)}.
\]
This differs from the normalized complexity \((f_{03}-20)/(f_0+f_3-10)\) of \citet{Z02}, but has the same finite limits along sequences with \(f_0+f_3\to\infty\). Let \(\mathcal L_4\) be the set of finite such limits. Let \(\mathcal H_{\mathrm{pow}}\) be the harmonic degrees of normal periodic weighted-Delaunay mosaics in \(\R^3\) arising from finitely many weighted sites per period; the same set is obtained from their Laguerre duals.

\begin{lemma}[General two-stack cost]\label{lem:generalport}
Let \(P\) have counts \((a,b,j)=(f_0,f_3,f_{03})\). Choose a facet \(G\) with \(v\) vertices and \(f\) polygonal faces, and a \(k\)-gonal face of \(G\). Stack beyond \(G\), then beyond the new pyramid on this \(k\)-gon. The resulting four-polytope has a tetrahedral facet and counts
\[
 a'=a+2,\qquad b'=b+f+k-1,\qquad
 j'=j+v+3f+4k-4.
\]
The inserted mosaic has
\begin{equation}\label{eq:generalport}
 h=\frac{j+v+3f+4k-8}{a+b+f+k-23/6}.
\end{equation}
\end{lemma}
\begin{proof}
If \(G\) has \(e\) edges, its face sizes sum to \(2e\). The first stacking changes the incidence count by \(2e+f-v\). Euler's relation \(v-e+f=2\) rewrites this as \(v+3f-4\). A pyramid on a \(k\)-gon has \(k+1\) vertices and \(k+1\) faces, namely the base and \(k\) triangles. Stacking it replaces it by one \((k+1)\)-vertex pyramid and \(k\) tetrahedra, adding \(4k\) incidences and \(k\) facets. There is at least one tetrahedral facet. The remaining counts and \eqref{eq:generalport} follow from Corollary~\ref{cor:polytocount}.
\end{proof}

\begin{lemma}[Boundary accounting for the capped epigraph]\label{lem:cap}
Fix a normal periodic weighted-Delaunay mosaic with finitely many cell types, convex lifting $H$, and densities $(V,C,I)$. For generic translated cubes $B_R$ of side $2R$, and $T_R>\max_{B_R}H$, the capped epigraph $P_R$ satisfies
\[
(f_0(P_R),f_3(P_R),f_{03}(P_R))=(V,C,I)\operatorname{vol}(B_R)+O(R^2)
\]
componentwise. The constants may depend on the fixed mosaic.
\end{lemma}
\begin{proof}
Choose a common cell-diameter bound $D$. Only cells intersecting the width-$D$ neighborhood of $\partial B_R$ can be clipped. Periodicity and the finite list of cell types bound their number by $O(R^2)$ and bound each cell's numbers of faces and incidences uniformly. Intersection with six halfspaces has bounded combinatorial complexity per such cell: each new vertex is determined by a face of that cell and a subset of the six clipping planes. Thus all new lower vertices and their incidences with clipped lower facets contribute $O(R^2)$. Interior vertices, facets and incidences are unaltered. Counts in the cube differ from periodic density times its volume by $O(R^2)$, by covering with fundamental domains and counting those meeting its boundary.

There are at most six vertical facets. Their lower vertices are among the just-counted boundary vertices; each is incident with at most six of these facets. The top facet has the eight lifted cube corners, since $T_R$ is strictly above the lower graph. The vertical facets have no other vertices away from their lower boundary and those top corners. Hence their additional vertices and vertex--facet incidences are $O(R^2)$, not a volume-order term. The seven added facets themselves contribute $O(1)$. These observations prove all three estimates, including the incidence estimate.
\end{proof}

\begin{theorem}[Asymptotic comparison]\label{thm:equivalence}
The finite closures satisfy
\[
 \overline{\mathcal H_{\mathrm{pow}}}=\mathcal L_4.
\]
Moreover, harmonic degree is unbounded on periodic weighted-Delaunay/Laguerre mosaics if and only if \(\rho(P)\) is unbounded on convex four-polytopes.
\end{theorem}
\begin{proof}
First let \(P_\ell\) satisfy \(D_\ell=a_\ell+b_\ell\to\infty\) and \(j_\ell/D_\ell\to H<\infty\). Choose a facet with the fewest vertices, so \(v_\ell\le j_\ell/b_\ell\). Every three-polytope satisfies \(f\le2v-4\), and has a polygonal face with at most five sides: \(e\le3f-6\) implies average face size \(2e/f<6\). Thus choose \(k\le5\).

Necessarily \(b_\ell\to\infty\): with a bounded number of facets, every vertex is determined by at least one independent four-tuple of facet hyperplanes, so the number of vertices is bounded by \(\binom{b_\ell}{4}\). It follows that
\[
 \frac{v_\ell}{D_\ell}\le\frac{j_\ell/D_\ell}{b_\ell}\longrightarrow0,
 \qquad \frac{f_\ell}{D_\ell}\longrightarrow0.
\]
The finite port cost in \eqref{eq:generalport} is therefore negligible. Regular implantation gives harmonic degrees tending to \(H\). This proves \(\mathcal L_4\subseteq\overline{\mathcal H_{\mathrm{pow}}}\).

Conversely, let a periodic weighted-Delaunay mosaic have density triple \((V,C,I)\), and let \(H(x)\) be its convex lower-hull lifting. It is quadratically periodic and has finitely many cell types. Restrict its epigraph to a generic large cube \(B_R\), and cap it above by a horizontal plane at height \(T_R>\max_{B_R}H\). This gives a convex four-polytope
\[
 P_R=\{(x,t):x\in B_R,\ H(x)\le t\le T_R\}.
\]
There are \(V\operatorname{vol}(B_R)+O(R^2)\) lower vertices and \(C\operatorname{vol}(B_R)+O(R^2)\) lower facets. Clipping affects only cells meeting the boundary layer, whose number and total incidence count are \(O(R^2)\). The six vertical side facets and the top facet add only \(O(R^2)\) incidences and vertices. Hence
\[
 f_0(P_R)=V\operatorname{vol}(B_R)+O(R^2),\quad
 f_3(P_R)=C\operatorname{vol}(B_R)+O(R^2),\quad
 f_{03}(P_R)=I\operatorname{vol}(B_R)+O(R^2).
\]
Lemma~\ref{lem:cap} supplies the full boundary-incidence accounting. Thus \(\rho(P_R)\to I/(V+C)\). To check closedness explicitly, let $H_n\in\mathcal L_4$ tend to a finite $H$. From a sequence defining $H_n$, choose $P_n$ with $f_0(P_n)+f_3(P_n)\ge n$ and $|\rho(P_n)-H_n|<1/n$. Then $\rho(P_n)\to H$, so $H\in\mathcal L_4$. This proves the reverse inclusion for closures.

If \(j/D\to\infty\), the same minimum-facet construction gives a harmonic degree bounded below by
\[
 \frac{j}{D+2j/b+5}.
\]
Its reciprocal is at most \(D/j+2/b+5/j\), which tends to zero. Thus unbounded polytope complexity implies unbounded harmonic degree. The capped-epigraph construction proves the converse. Bounded \(D\) allows only bounded flag counts, so no unbounded sequence is lost by imposing \(D\to\infty\).
\end{proof}

\begin{corollary}[Lower bound in the periodic regular class]\label{cor:lower}
Every normal periodic weighted-Delaunay mosaic in the class defining \(\mathcal H_{\mathrm{pow}}\), and every corresponding Laguerre dual, has \(\hbarM\ge3\).
\end{corollary}
\begin{proof}
The established four-polytope inequality
\[
 f_{03}-3(f_0+f_3)\ge-10
\]
is the nonnegativity of the second toric \(g\)-number; see \citet[Section 4]{Z02}, with the general convex case attributed there to Kalai's rigidity argument. Apply it to the capped epigraph polytopes and divide by \(\operatorname{vol}(B_R)\). Taking the limit gives \(I\ge3(V+C)\). Duality preserves the harmonic degree.
\end{proof}

This corollary is restricted to the stated liftable periodic class. It is not a proof of the conjectured strict lower bound for arbitrary convex mosaics. Theorem~\ref{thm:equivalence} explains precisely what a complete upper-range classification in the periodic power-diagram class would require. It does not supply a numerical maximum or decide whether the range is unbounded.

\section{Verification and interpretation}\label{sec:verification}
The construction and its verification use exact rational arithmetic. The finite certificate does not rely on numerical convex-hull tolerances, a picture of a tiling, or incidence feasibility alone.

\begin{center}
\begin{tabular}{p{0.52\textwidth}p{0.37\textwidth}}
\toprule
Check & Exact result\\
\midrule
Projected deformed-cube positive-circuit rays &66 total; 64 supporting facets; 2 strict redundant inequalities\\
Initial four-polytope face vector &\((64,192,192,64)\)\\
Twice-stacked face vector &\((66,205,212,73)\)\\
Template faces &208 internal polygons; 4 boundary triangles\\
Template cells and incidences &72 and 546\\
Sum of exact template cell volumes &\(1/6\)\\
Template regular lifting &Every cell is an exact strict lower facet\\
Periodic weighted realization &373 sites, 432 Delaunay cells; 24 positive macro-face tests\\
Harmonic degree &\(468/115\)\\
\bottomrule
\end{tabular}
\end{center}

The verifier reconstructs all nonnegative elimination circuits of support at most three, checks every facet's full equality set and every vertex's extremality, enumerates the face lattice by intersections, replays the stackings, and checks the projection in exact barycentric coordinates. It then checks interior-face multiplicities and triangulates every polygonal cell boundary through face and cell barycenters to obtain exact volumes. The weighted realization is checked independently through its local convexity inequalities.

\paragraph{Why nonsimplicial cells matter.}
For any periodic tetrahedral mosaic, $I=4C$, hence $\hbarM=4C/(V+C)<4$. Arbitrarily triangulating our nonsimplicial cells therefore cannot preserve the counterexample's harmonic degree. The nonsimplicial incidence structure is substantive, not a drawing convention.

\paragraph{Geological scope.}
A universal statement about all normal convex face-to-face mosaics is false. This does not say that the constructed highly structured cells are typical rock fragments, that an ordinary fracture mechanism generates them, or that empirical observations of a narrow band are incorrect. Normality is a geometric size condition, not a stochastic fracture law or a near-isotropy assumption. A physically motivated replacement theorem must add explicit restrictions and test them; convexity and face matching alone do not imply the proposed band.

\paragraph{Claim boundary.}
The finite upper-band counterexample, periodic Laguerre realization, explicit unbounded-degree families, interval inclusion \((3,16)\), and regular-class asymptotic comparison are proved above. We do not claim an exact description of \(S_3\), a universal maximum of 16, the strict lower bound \(\hbarM>3\) for all convex mosaics, or a counterexample within unweighted Euclidean Voronoi mosaics. The code is a reproducible exact check, not a proof-assistant formalization or an independent human referee report.

\appendix
\section{A self-contained rational four-polytope certificate}\label{sec:rational}
Set \(\varepsilon=1/2\). In \(\R^6\) impose
\begin{equation}\label{eq:deformed}
 \varepsilon |x_k|\le
 \frac{2^{\binom{k}{2}}}{\varepsilon^{k-1}}
 -(-1)^k\sum_{j=1}^{k-1}\binom{k-2}{j-1}x_j,
 \qquad 1\le k\le6,
\end{equation}
where the sum is empty for \(k=1\). This is a specialization of the deformed-cube construction in \citet{JZ}; we verify the specialization rather than relying on an unspecified ``sufficiently small'' parameter.

For each sign vector \(\sigma\in\{-1,1\}^6\), solve the equalities recursively, selecting the sign \(\sigma_k\) at step \(k\). At every prefix vertex the right side is strictly positive. Since it is affine in the prefix variables, it remains positive throughout the prefix polytope. Induction therefore shows that the feasible set is a combinatorial six-cube and that the 64 recursively generated points are all its vertices. Project to the last four coordinates and rescale them by
\[
 (x_3,x_4,x_5,x_6)\mapsto
 \left(\frac{x_3}{64},\frac{x_4}{1024},\frac{x_5}{32768},\frac{x_6}{2097152}\right).
\]
The result is the initial rational four-polytope used in the certificate.

\subsection{Why the facet enumeration is complete}
There are twelve signed inequalities in \eqref{eq:deformed}. Their first two coefficients, the eliminated coordinates, are
\[
 (\pm\tfrac12,0),\quad(1,\pm\tfrac12),\quad
 (-1,-1),\quad(1,2),\quad(-1,-3),\quad(1,4),
\]
with each of the last four vectors appearing twice. The nonnegative dependence cone of these vectors consists of nonnegative multipliers which cancel the first two coordinates. Its extreme rays have support at most three: a minimal positive dependence in a two-dimensional vector configuration has at most three members. Enumerating all subsets of sizes one, two and three therefore enumerates all extreme elimination rays.

For each ray, combine its original inequalities. Exact evaluation on the 64 projected vertices finds 66 rays: two resulting inequalities are strict everywhere on the projected polytope, while each of the other 64 has exactly eight equality vertices spanning a three-dimensional affine space. These 64 equality sets are distinct. Every projected facet is exposed by an extreme elimination ray or a sum of rays exposing the same facet, so the list is complete. This is also an immediate application of the linear-inequality alternative underlying projection/elimination: the projected set is described by all nonnegative combinations with vanishing eliminated coordinates.

The verifier checks that the normals incident at every listed vertex span four dimensions. It computes all face sets as intersections of the facet vertex sets and obtains \((64,192,192,64)\). Within each facet it obtains eight vertices, twelve edges and six quadrilateral faces; all degrees are those of a cube. Thus no combinatorial or realizability assumption is left unchecked in the finite example.

\subsection{Deterministic stacking and projection}
Facets are ordered lexicographically by their vertex-index sets. Stack first beyond the first facet. For any target facet, let \(z\) be its vertex barycenter and \(w\) the mean of all current polytope vertices. Choose the first number in \(1/2,1/4,1/8,\ldots\) for which
\[
 p=z+\delta(z-w)
\]
is strictly beneath every other facet. Since \(w\) is interior and \(z\) is in the target facet's relative interior, such a choice exists and is beyond the target facet. The verifier checks every sign explicitly.

The second target is the first square-pyramid facet containing the first new vertex. The outer facet for the Schlegel projection is the first tetrahedral facet containing the second new vertex. The viewpoint is chosen by the same dyadic rule. Its exact projection formula is given in Lemma~\ref{lem:schlegel}. An affine barycentric coordinate map sends the four outer vertices, in their stored order, to \(0,e_1,e_2,e_3\).

The file \texttt{mosaic\_certificate.json} contains all 64 initial vertices and supporting halfspaces, elimination circuits, both stacking records, all 66 final four-dimensional vertices, all 73 facets, the viewpoint, the 66 three-dimensional template coordinates and their regular heights. Thus the finite object can be regenerated without network access or a polytope library.

\section{Reproduction instructions}
Run from the package root, using Python with assertions enabled:
\begin{verbatim}
python code/build_mosaic.py
python code/verify_mosaic.py
python code/verify_regular_realization.py
python code/check_weighted_certificate.py
python code/verify_abrasion.py
python code/test_rejection.py
\end{verbatim}
The only nonstandard runtime dependency is SymPy. The scripts reject optimized Python execution, which would disable assertions. The final three JSON reports and the rejection-test report are in \texttt{evidence/}. The abrasion verifier belongs to the separate companion note. The package manifest records SHA-256 hashes of the delivered files. Mathematical assertions about infinite families and PDE evolution are established in the manuscripts; the scripts certify the finite identities on which their stated arguments rely.

\paragraph{Construction versus checking.}
The historically named \path{verify_regular_realization.py} is a checked constructor: it regenerates and overwrites the weighted-site payload. The separate \path{check_weighted_certificate.py} is read-only and compares the entire supplied payload to an exact checked reconstruction. Rejection controls alter weights, lattice translations and cell references, and verify input-byte preservation. Both routes share reconstruction code and SymPy; this is not an independent implementation of the infinite lower-hull proof.

\begin{thebibliography}{9}
\bibitem[Rybnikov(1999)]{R99}
Rybnikov, K. (1999). Stresses and liftings of cell-complexes. \emph{Discrete \& Computational Geometry, 21}, 481--517. \url{https://doi.org/10.1007/PL00009434}.
\bibitem[Domokos and L\'angi(2022)]{DL}
Domokos, G., \& L\'angi, Z. (2022). On some average properties of convex mosaics. \emph{Experimental Mathematics, 31}(3), 783--793. Published online 2019. \url{https://doi.org/10.1080/10586458.2019.1691090}. Author version: \url{https://arxiv.org/abs/1905.00721}.

\bibitem[Joswig and Ziegler(2000)]{JZ}
Joswig, M., \& Ziegler, G. M. (2000). Neighborly cubical polytopes. \emph{Discrete \& Computational Geometry, 24}, 325--344. \url{https://doi.org/10.1007/s004540010039}. \url{https://arxiv.org/abs/math/9812033}.

\bibitem[Ziegler(2002)]{Z02}
Ziegler, G. M. (2002). Face numbers of 4-polytopes and 3-spheres. In \emph{Proceedings of the International Congress of Mathematicians, Beijing 2002} (Vol. III, pp. 625--634). Higher Education Press. Revised author version (2003): \url{https://arxiv.org/abs/math/0208073}.

\bibitem[Ziegler(2004)]{Z04}
Ziegler, G. M. (2004). Projected products of polygons. \emph{Electronic Research Announcements of the American Mathematical Society, 10}, 122--134. \url{https://doi.org/10.1090/S1079-6762-04-00137-4}. Extended abstract: \url{https://arxiv.org/abs/math/0407042} (arXiv metadata title: \emph{Projected Products of Polytopes}).
\end{thebibliography}
\end{document}
